\section{Исследование и построение решения задачи}
\label{sec:Chapter3} \index{Chapter3}

%Требуется разбить большую задачу, описанную в постановке, на более мелкие
%подзадачи. Процесс декомпозиции следует продолжать до тех пор, пока подзадачи
%не станут достаточно простыми для решения непосредственно. Это может быть
%достигнуто, например, путем проведения эксперимента, доказательства теоремы
%или поиска готового решения.

Пусть $\mu, \nu \sim \chi_k$ --- независимые случайные величины в $\Z_q$. Обозначим распределение их произведения $\mu \cdot \nu$ за $\xi_k$. Положим $\psi = \xi_k^{\circledast n} * \xi_k^{\circledast n} * \chi_k$. В~статье~\cite{9457596} доказывалось, что это распределение коэффициентов $\bold z$, $\reqv$
\begin{equation*}
	\Pr(\mathcal{E}) \leq 1-\sum_{i=-b}^{b}{\psi[i]} \,\;\text{для}\;\, b = \left\lfloor\frac{q}{2Q}\right\rfloor; \quad \mathsf{DFR} \leq \sum_{i=t+1}^n{C_n^i\Pr(\mathcal{E})^i(1-\Pr(\mathcal{E}))^{n-i}}
\end{equation*}
Убрать индексы у $\mathcal{E}_i$ корректно, так как распределение у них одинаковое. При этом во втором неравенстве предполагается их независимость. Математически это, конечно, не так, однако на практике близко к оному.
